\clearpage
\VerdierPhysicalPageAnchor{0260}%
\begingroup
\small
\setlength{\parindent}{0pt}%
\setlength{\parskip}{0pt}%
\setlength{\emergencystretch}{2em}%
\newcommand{\VerdierTermRow}[3]{%
  \VerdierSemanticAnchor{#1}%
  \par\noindent\hangindent=1.5em\hangafter=1
  #2。\enspace\textup{#3。}\par\smallskip}

\VerdierTermRow{idxt-p260-r01}{アーベル圏のコホモロジー次元}{第 III 章、3.1.1}
\VerdierTermRow{idxt-p260-r02}{複体の圏への関手の拡張}{第 I 章、2.4.2}
\VerdierTermRow{idxt-p260-r03}{スペクトル対象に付随するスペクトル族}{第 II 章、4.2.5}
\VerdierTermRow{idxt-p260-r04}{第 $m$ 次で定まるフィルトレーション}{第 III 章、4.1.6}
\VerdierTermRow{idxt-p260-r05}{反完全関手}{第 II 章、1.1.6}
\VerdierTermRow{idxt-p260-r06}{コホモロジー的関手}{第 II 章、1.1.5}
\VerdierTermRow{idxt-p260-r07}{標準コホモロジー的関手}{第 II 章、1.3.5；第 III 章、1.2.2}
\VerdierTermRow{idxt-p260-r08}{二つの $G$-関手の合成}{第 I 章、1.2.6}
\VerdierTermRow{idxt-p260-r09}{写像錐関手}{第 I 章、3.1.2}
\VerdierTermRow{idxt-p260-r10}{群の変更関手}{第 I 章、1.4.1}
\VerdierTermRow{idxt-p260-r11}{局所化関手}{第 II 章、2.2.5}
\VerdierTermRow{idxt-p260-r12}{商関手}{第 II 章、2.2.10}
\VerdierTermRow{idxt-p260-r13}{方向 $i$ のシフト関手}{第 I 章、2.3.1}
\VerdierTermRow{idxt-p260-r14}{次数関手}{第 I 章、1.1.1}
\VerdierTermRow{idxt-p260-r15}{完全関手}{第 II 章、1.1.3}
\VerdierTermRow{idxt-p260-r16}{定常完全（それぞれコホモロジー的）関手（それぞれ左定常、右定常）}{第 III 章、4.5.1}
\VerdierTermRow{idxt-p260-r17}{次数付き関手}{第 I 章、1.1.1}
\VerdierTermRow{idxt-p260-r18}{加法的次数付き関手}{第 I 章、1.1.3}
\VerdierTermRow{idxt-p260-r19}{捩れ $G$-関手の反対捩れ次数付き関手}{第 I 章、1.4.8}
\VerdierTermRow{idxt-p260-r20}{多重完全関手}{第 II 章、1.1.4}
\VerdierTermRow{idxt-p260-r21}{$G$-圏}{第 I 章、1.2.2}
\VerdierTermRow{idxt-p260-r22}{加法的 $G$-圏}{第 I 章、1.2.8}
\VerdierTermRow{idxt-p260-r23}{厳密な $G$-圏}{第 I 章、1.2.2}
\VerdierTermRow{idxt-p260-r24}{$G$-関手}{第 I 章、1.2.4}
\VerdierTermRow{idxt-p260-r25}{厳密な $G$-関手}{第 I 章、1.2.4}
\VerdierTermRow{idxt-p260-r26}{余サイクルで捩られた $G$-関手}{第 I 章、1.4.4}
\VerdierTermRow{idxt-p260-r27}{二つのホモトピーを結ぶ二階ホモトピー}{第 I 章、2.5.12}
\endgroup
